\chapter{Introducción}\label{chapter:introduccion}

La industria de los videojuegos ha experimentado un crecimiento sostenido durante las últimas décadas, generando una oferta cada vez más amplia y diversa de títulos disponibles para los jugadores. Si bien esta gran variedad representa una ventaja para los usuarios, también genera dificultades a la hora de decidir qué videojuego jugar, debido a la sobrecarga de información y la cantidad de alternativas existentes.

En este contexto, los sistemas de recomendación se han convertido en herramientas fundamentales para ayudar a los usuarios a descubrir contenido alineado con sus gustos e intereses. Sin embargo, muchas de las soluciones actuales presentan limitaciones relacionadas con la personalización, la explicabilidad de las recomendaciones y la construcción de perfiles detallados de los jugadores.

Con el objetivo de abordar esta problemática, el presente proyecto propone el desarrollo de GameTrack, una plataforma orientada al registro, análisis y recomendación personalizada de videojuegos mediante el uso de técnicas de perfilado de usuarios, sistemas de recomendación e inteligencia artificial. La solución busca asistir a los jugadores en la toma de decisiones, facilitar el descubrimiento de nuevos títulos y ofrecer herramientas de análisis que aporten valor tanto a usuarios como a desarrolladores de videojuegos.

\section{Objetivos}\label{sec:objetivos}

Son objetivos generales de este proyecto, de ahora en más, GameTrack:

• Desarrollar una plataforma capaz de asistir a los usuarios en la elección de videojuegos mediante recomendaciones personalizadas.

• Implementar un sistema de recomendación basado en técnicas de perfilado de usuarios, inteligencia artificial y análisis de preferencias.

• Facilitar el descubrimiento de videojuegos alineados con los intereses y características de cada jugador.

Son objetivos específicos de GameTrack:

• Construir perfiles de usuario a partir de información explícita e implícita obtenida mediante interacciones, calificaciones y reseñas.

• Diseñar e implementar un sistema de recomendación híbrido que combine diferentes técnicas de recomendación para mejorar la precisión de las sugerencias.

• Desarrollar un asistente inteligente que ayude a los usuarios a decidir qué videojuego jugar según sus preferencias y contexto actual.

• Implementar funcionalidades de recomendación grupal que permitan identificar videojuegos compatibles con los intereses de múltiples usuarios.

• Proporcionar herramientas de análisis y segmentación que permitan a los desarrolladores comprender mejor a sus potenciales audiencias.

• Evaluar el desempeño del sistema mediante métricas de recomendación y pruebas sobre usuarios potenciales.

\section{Alcance}\label{sec:alcance}

El presente proyecto contempla el diseño y desarrollo de un prototipo funcional denominado GameTrack, orientado a la recomendación personalizada de videojuegos mediante técnicas de perfilado de usuarios, sistemas de recomendación e inteligencia artificial.

La solución incluirá funcionalidades de registro e inicio de sesión de usuarios, construcción de perfiles mediante información explícita e implícita, gestión de videojuegos, calificaciones, reseñas y listas personales. Asimismo, se desarrollará un sistema de recomendación híbrido capaz de generar sugerencias personalizadas a partir de las preferencias e interacciones de los usuarios, incorporando además un asistente inteligente para facilitar la decisión de qué videojuego jugar y funcionalidades de recomendación grupal para múltiples usuarios.

El proyecto también contempla herramientas de análisis orientadas a desarrolladores, permitiendo obtener información sobre perfiles de jugadores, preferencias generales y tendencias identificadas a partir de los datos recopilados por la plataforma.

Como resultado, se espera obtener un prototipo funcional que permita validar la viabilidad técnica y conceptual de la propuesta. Quedan fuera del alcance del proyecto aspectos relacionados con el despliegue comercial a gran escala, la integración con múltiples plataformas externas más allá de las necesarias para la validación del sistema, la optimización para una gran masa de usuarios concurrentes y el desarrollo de modelos de inteligencia artificial de nivel industrial.
